% !TEX program = xelatex
\documentclass[UTF8,zihao=-4]{ctexart}

\usepackage[a4paper,margin=2.3cm]{geometry}
\usepackage{amsmath,amssymb}
\usepackage{booktabs,longtable,array}
\usepackage{xcolor}
\usepackage{hyperref}
\usepackage{listings}
\usepackage{fancyhdr}
\usepackage{enumitem}

\definecolor{codebg}{RGB}{247,248,250}
\definecolor{codeframe}{RGB}{213,217,223}
\definecolor{accent}{RGB}{30,83,150}
\definecolor{warn}{RGB}{165,95,20}

\newcolumntype{L}[1]{>{\raggedright\arraybackslash}p{#1}}
\setlist{nosep,leftmargin=2em}

\lstdefinestyle{code}{
  basicstyle=\ttfamily\small,
  backgroundcolor=\color{codebg},
  frame=single,
  rulecolor=\color{codeframe},
  breaklines=true,
  columns=fullflexible,
  keepspaces=true,
  showstringspaces=false,
  tabsize=2
}

\hypersetup{
  colorlinks=true,
  linkcolor=accent,
  urlcolor=accent,
  citecolor=accent,
  pdftitle={EviMind 项目解析教程},
  pdfauthor={Codex}
}

\setlength{\headheight}{15pt}
\pagestyle{fancy}
\fancyhf{}
\lhead{EviMind 项目解析教程}
\rhead{\thepage}

\newcommand{\term}[1]{\textbf{#1}}
\newcommand{\risk}[1]{\textcolor{warn}{\textbf{#1}}}

\title{\textbf{EviMind 项目解析教程}\\
\large 从启动、架构、核心链路到二次开发}
\author{}
\date{\today}

\begin{document}
\maketitle

\begin{abstract}
这是一份面向初学者的 EviMind 项目解析教程。读者不需要先熟悉 Spring Boot、Vue、RAG 和向量数据库的全部细节，只要按顺序阅读，就能理解项目定位、运行方式、目录结构、后端主链路、前端请求层、数据库模型、文档入库、RAG 问答、安全控制、部署方式和验证边界。本文基于当前仓库文件写成，避免脱离代码的泛泛介绍。
\end{abstract}

\tableofcontents
\newpage

\section{先看项目定位}

EviMind 是一个证据驱动的 RAG 知识工作台。它的核心目标不是单纯聊天，而是把文档上传、文本抽取、切片、检索、证据引用、研究笔记、引用导出和批量分析放进同一个工作流。

一句话理解：

\begin{quote}
用户上传资料，系统把资料变成可检索证据；用户提问时，系统先检索证据，再让大模型基于证据回答，并返回引用。
\end{quote}

当前项目适合三类场景：

\begin{itemize}
  \item 本地研究资料管理：上传 PDF、Word、Markdown、代码、日志等文件，围绕知识库问答。
  \item RAG 工程学习：能看到文件入库、混合检索、RRF 融合、证据门控、SSE 流式输出的完整实现。
  \item 项目交付演示：支持 standalone、单 JAR 和 Docker Compose 三种运行模式。
\end{itemize}

\section{新手阅读路线}

不要从所有类开始看。这个项目模块多，直接翻目录容易乱。建议按下面顺序：

\begin{enumerate}
  \item 先跑 standalone，确认登录、知识库、上传、问答四件事能串起来。
  \item 看 README 的 Core Workflow，理解业务主线。
  \item 看认证链路：JWT 过滤器、GroupContext、知识库成员校验。
  \item 看文档上传链路：DocumentController 到 DocumentService，再到 EtlPipeline。
  \item 看问答链路：ConversationController 到 ConversationService，再到 RagPipeline。
  \item 看检索链路：HybridSearchService、PgVectorSearchService、ElasticsearchSearchService、RrfFusionService。
  \item 最后看前端：路由、请求封装、ChatView、各模块 API 文件。
\end{enumerate}

这样读的好处是先抓住数据流，再回头补框架细节。

\section{运行方式}

\subsection{开发模式}

开发模式使用 backend standalone profile，数据库走 H2，文件走本地目录。PostgreSQL、Elasticsearch、MinIO 都不是必需。

\begin{lstlisting}[style=code]
cd D:\EviMind
$env:JAVA_HOME = "C:\Program Files\Java\jdk-22"
$env:DEEPSEEK_API_KEY = "your-key"
.\mvnw.cmd spring-boot:run "-Dspring-boot.run.profiles=standalone"
\end{lstlisting}

前端单独启动：

\begin{lstlisting}[style=code]
cd D:\EviMind\frontend
npm install
npm run dev -- --host 127.0.0.1
\end{lstlisting}

浏览器打开：

\begin{lstlisting}[style=code]
http://127.0.0.1:5173
\end{lstlisting}

前端 Vite 会把 \path{/api} 代理到 \path{http://localhost:8080}。后端健康检查是：

\begin{lstlisting}[style=code]
http://localhost:8080/actuator/health
\end{lstlisting}

\subsection{单 JAR 模式}

单 JAR 模式适合本地交付。脚本会先构建前端，把 \path{frontend/dist} 复制到 \path{src/main/resources/static}，再打 Spring Boot JAR。

\begin{lstlisting}[style=code]
cd D:\EviMind
$env:JAVA_HOME = "C:\Program Files\Java\jdk-22"
$env:DEEPSEEK_API_KEY = "your-key"
.\build.bat
.\start.bat
\end{lstlisting}

访问入口：

\begin{lstlisting}[style=code]
http://localhost:8080
\end{lstlisting}

\subsection{Docker Compose 模式}

Docker Compose 模式启动 PostgreSQL/pgvector、Elasticsearch、MinIO、backend 和 frontend。它比 standalone 更接近完整部署环境。

\begin{lstlisting}[style=code]
cd D:\EviMind
copy .env.example .env
docker compose up -d
\end{lstlisting}

启动前必须在 \path{.env} 填入数据库、MinIO 和 JWT 密钥。Compose 缺少这些值会拒绝解析；后端和基础设施端口仅绑定到本机回环地址。

\section{技术栈总览}

\begin{longtable}{L{0.22\linewidth}L{0.68\linewidth}}
\toprule
层次 & 当前实现 \\
\midrule
后端框架 & Spring Boot 3.5，Java 21，Spring Web，Spring Security，Actuator \\
AI 接入 & Spring AI 1.0.0-M1，OpenAI-compatible ChatClient，多 provider 配置 \\
数据访问 & MyBatis-Plus 为主，部分 JPA Repository 用于分析结果 \\
数据库 & PostgreSQL/pgvector 用于完整栈，H2 用于 standalone \\
搜索 & pgvector 语义检索、Elasticsearch 关键词检索、本地关键词回退 \\
对象存储 & MinIO 或本地文件存储 \\
文档解析 & PDFBox、Apache POI、OpenPDF、Tess4J、epub4j \\
前端 & Vue 3，TypeScript，Vite，Element Plus，Pinia，Axios，Markdown-It \\
工程质量 & Maven Surefire/Failsafe、JaCoCo、Spotless、PMD、GitHub Actions \\
\bottomrule
\end{longtable}

初学者可以先不深究每个依赖。先记住：Spring Boot 负责 API，Vue 负责界面，PostgreSQL/pgvector 与 Elasticsearch 负责检索，MinIO/本地目录负责文件，Spring AI 负责模型调用。

\section{目录结构}

仓库主要目录如下：

\begin{lstlisting}[style=code]
EviMind/
  src/main/java/com/example/evimind/
    assistant/       会话、流式对话、RAG 问答入口
    auth/            注册、登录、JWT、刷新令牌
    common/          健康检查、全局异常处理
    config/          AI、安全、异步、存储、搜索配置
    document/        文档上传、列表、删除、重试
    extractor/       PDF、Word、Excel、PPT、EPUB、图片文字提取
    ingestion/       ETL 入库流水线、切片、嵌入、ES 写入
    knowledgebase/   知识库 CRUD、成员、检索入口
    qa/              RAG 编排、证据组合、回答与引用
    retrieval/       查询改写、向量检索、关键词检索、RRF
    service/         分析、引用、报告、知识图谱等业务服务
    storage/         MinIO 和本地文件存储
  src/main/resources/
    application.yml
    application-standalone.yml
    schema-h2.sql
    db/migration/
    prompts/
  frontend/
    src/views/
    src/stores/
    src/api/
    src/utils/
\end{lstlisting}

读目录时注意两个点：

\begin{itemize}
  \item \path{controller} 和 \path{assistant/document/knowledgebase} 里都有 REST 入口，不能只看一个包。
  \item \path{schema-h2.sql} 服务 standalone；\path{db/migration} 服务 PostgreSQL/Flyway。两者覆盖范围不同。
\end{itemize}

\section{后端启动入口}

后端入口是 \path{EvimindApplication.java}。当前类排除了 Spring AI 的部分自动配置：

\begin{lstlisting}[style=code]
@SpringBootApplication(exclude = {
  OpenAiAutoConfiguration.class,
  ChatClientAutoConfiguration.class
})
\end{lstlisting}

这表示项目自己创建 ChatClient，而不是完全依赖默认自动装配。具体实现看 \path{config/AiConfig.java}：

\begin{itemize}
  \item 读取 \path{custom.ai.providers} 下的 deepseek、zhipu、qianwen、openai 配置。
  \item 为每个 provider 创建 OpenAI-compatible \path{OpenAiApi}。
  \item 生成 \path{Map<String, ChatClient>}，业务层按 provider 取模型。
  \item 如果开启 embedding，再创建 \path{OpenAiEmbeddingModel}。
\end{itemize}

\risk{当前检查点：} \path{AiDataTools.java} 定义了 listDirectory、readFileContent、queryAnalysisHistory 三个函数 Bean，但当前 \path{AiConfig.java} 没有看到把这些函数显式绑定到 ChatClient 的调用。以后如果要做真正 agent 工具调用，需要补 ChatClient 函数绑定，并加测试验证。

\section{配置分层}

\subsection{默认配置}

\path{application.yml} 面向完整栈。它默认使用 PostgreSQL、Flyway、MinIO、Elasticsearch，embedding 默认开启，JWT secret 通过环境变量注入。

关键配置分组：

\begin{longtable}{L{0.28\linewidth}L{0.60\linewidth}}
\toprule
配置路径 & 作用 \\
\midrule
\path{spring.datasource} & PostgreSQL 连接 \\
\path{spring.flyway} & 数据库迁移 \\
\path{spring.elasticsearch.uris} & Elasticsearch 地址 \\
\path{minio} & 对象存储 \\
\path{jwt} & access token 和 refresh token 生命周期 \\
\path{custom.ai.providers} & DeepSeek、GLM、Qianwen、OpenAI-style provider \\
\path{custom.ai.embedding} & 向量模型配置 \\
\path{custom.rag} & 检索超时、查询改写、rerank、证据上下文长度 \\
\path{custom.analysis} & 批量分析和报告导出目录 \\
\bottomrule
\end{longtable}

\subsection{Standalone 配置}

\path{application-standalone.yml} 面向本地演示：

\begin{itemize}
  \item 数据库：\path{data/evimind-standalone.mv.db}。
  \item 文件目录：\path{data/documents}。
  \item MinIO：关闭。
  \item Elasticsearch：默认关闭。
  \item Embedding：默认关闭。
  \item H2 Console：开启，路径是 \path{/h2-console}。
\end{itemize}

这能降低启动成本，但也带来边界：H2 schema 当前只覆盖核心表；完整 PostgreSQL 迁移还包含向量、引用、知识图谱、权限、审计、定时任务等表。做完整能力验证时应使用 Docker/production-like 栈。

\section{认证和权限主线}

\subsection{注册和登录}

认证入口在 \path{AuthController.java}，核心逻辑在 \path{AuthService.java}。

\begin{enumerate}
  \item register：检查用户名唯一，BCrypt 加密密码，创建用户，返回 access token 和 refresh token。
  \item login：校验用户名密码，禁用用户拒绝登录，返回新 token。
  \item refresh：用 refresh token hash 查表，旧 refresh token 置为 revoked，再发新 token。
  \item me：从 GroupContext 读取当前用户 ID，再查用户信息。
\end{enumerate}

refresh token 不明文入库，入库的是 SHA-256 hash。这比直接保存 refresh token 更安全。

\subsection{JWT 过滤器}

\path{JwtAuthenticationFilter.java} 做三件事：

\begin{enumerate}
  \item 从 Authorization header 提取 Bearer token。
  \item 校验 token，取出 userId、username、systemRole。
  \item 调用 GroupService 获取默认 groupId，把 userId、groupId、role 写进 GroupContext。
\end{enumerate}

请求结束后，filter 在 finally 中清理 GroupContext。这个细节重要，因为 GroupContext 是 ThreadLocal，不清理会污染后续请求。

\subsection{业务权限}

后端权限不是只靠 SecurityConfig。很多业务服务会再次检查成员关系：

\begin{itemize}
  \item KnowledgeBaseService：创建者是 OWNER；更新、删除、成员管理要求 OWNER。
  \item DocumentService：上传、列表、详情、删除、重试前检查知识库成员。
  \item RagPipeline：问答前检查用户是否是知识库成员。
  \item ConversationService：会话只能由 owner 访问。
\end{itemize}

新手调试 403 时，不要只看 Spring Security。更常见的原因是 GroupContext 没有用户，或 kb\_member 没有对应记录。

\section{前端结构}

前端入口在 \path{frontend/src/main.ts}，路由在 \path{frontend/src/router/index.ts}。路由守卫逻辑很直接：需要登录的页面如果没有 accessToken，就跳到 \path{/login}。

主要页面：

\begin{longtable}{L{0.30\linewidth}L{0.58\linewidth}}
\toprule
页面 & 作用 \\
\midrule
\path{ChatView.vue} & 证据问答、模型参数、SSE 流式输出、引用展示 \\
\path{KnowledgeBaseView.vue} & 知识库管理 \\
\path{DocumentView.vue} & 文档上传、列表、入库状态、重试 \\
\path{AnalysisView.vue} & 文件分析和批量分析 \\
\path{CitationView.vue} & 引用导出 \\
\path{AcademicGraphView.vue} & 学术引用图 \\
\path{KnowledgeGraphView.vue} & 知识图谱 \\
\path{NotesView.vue} & 研究笔记 \\
\bottomrule
\end{longtable}

\subsection{请求封装}

\path{frontend/src/utils/request.ts} 是前端理解后端的关键文件：

\begin{itemize}
  \item Axios baseURL 是 \path{/api/v1}。
  \item 每个请求自动带 Authorization header。
  \item 遇到 401 且未重试过，会调用 \path{/api/v1/auth/refresh} 刷新 access token。
  \item SSE 不能完全靠 Axios，所以提供 \path{authenticatedFetch} 给流式请求使用。
\end{itemize}

这解释了为什么普通 API 用 \path{request.ts}，流式聊天用 fetch 读流。

\section{核心业务流程}

\subsection{总流程}

完整工作流如下：

\begin{lstlisting}[style=code]
register / login
  -> create knowledge base
  -> upload document
  -> async ETL
  -> create conversation
  -> send question
  -> hybrid retrieval
  -> evidence-gated generation
  -> stream answer + citations
\end{lstlisting}

这条线就是阅读项目的主干。

\subsection{知识库创建}

\path{KnowledgeBaseService.create} 会填默认值：

\begin{itemize}
  \item evidenceThreshold 默认 0.50。
  \item chunkStrategy 默认 PARAGRAPH。
  \item chunkSize 默认 500。
  \item chunkOverlap 默认 100。
  \item 创建者写入 kb\_member，角色是 OWNER。
\end{itemize}

知识库不是单纯分类目录，它还保存切片策略和证据阈值。也就是说，同一份文档放进不同知识库，后续入库和问答策略可能不同。

\section{文档上传和 ETL}

\subsection{上传入口}

文档上传入口是：

\begin{lstlisting}[style=code]
POST /api/v1/documents/upload
\end{lstlisting}

调用链：

\begin{lstlisting}[style=code]
DocumentController
  -> DocumentService.upload
  -> storage upload
  -> insert document(status=UPLOADED)
  -> triggerIngestionAsync
  -> EtlPipeline.processDocument
\end{lstlisting}

\path{DocumentService} 会先做安全检查：

\begin{itemize}
  \item 用户必须是知识库成员。
  \item 文件不能为空。
  \item 文件名会被清洗，路径穿越会被拒绝。
  \item 格式必须在白名单内。
  \item 文件不能超过 50MB。
  \item 常见扩展名和 MIME 类型要匹配。
\end{itemize}

存储路径不使用用户原始文件名，而是：

\begin{lstlisting}[style=code]
knowledgeBaseId/UUID/UUID.ext
\end{lstlisting}

原始文件名只作为展示元数据保存。

\subsection{ETL 状态机}

\path{EtlPipeline} 的状态顺序：

\begin{lstlisting}[style=code]
UPLOADED
  -> EXTRACTING
  -> CLEANING
  -> CHUNKING
  -> PERSISTING
  -> EMBEDDING
  -> INDEXING
  -> ENRICHING
  -> COMPLETED
\end{lstlisting}

失败时状态变为 FAILED，并记录 failedStage、errorCode、errorMessage、retryCount、startedAt、finishedAt、ingestionVersion。

\subsection{版本化入库}

项目使用 ingestionVersion 和 activeIngestionVersion 管理文档版本。重试入库时会产生新版本，只有新版本完成后才切换 activeIngestionVersion。检索时只读 active version，避免失败重试留下的旧 chunk 被召回。

这个设计的价值是：入库过程可以失败，但失败不应该污染线上可检索证据。

\subsection{切片策略}

\path{DocumentChunker.java} 支持三种策略：

\begin{longtable}{L{0.22\linewidth}L{0.65\linewidth}}
\toprule
策略 & 含义 \\
\midrule
FIXED\_LENGTH & 按固定长度切，简单但可能切断语义 \\
PARAGRAPH & 按段落组织，再受 chunkSize 和 overlap 约束 \\
SEMANTIC & 按句子聚合，更适合自然语言文本 \\
\bottomrule
\end{longtable}

overlap 不是越大越好。过大带来重复召回和成本上升；过小可能丢上下文。默认 500/100 是一个入门可用值。

\section{检索链路}

\subsection{HybridSearchService}

检索链路在 \path{HybridSearchService.java}：

\begin{lstlisting}[style=code]
user query
  -> QueryRewriteService
  -> pgvector semantic search
  -> Elasticsearch keyword search
  -> local keyword fallback
  -> RRF fusion
\end{lstlisting}

两路检索通过 CompletableFuture 并行执行，默认后端超时是 1500ms。任一路失败都不会直接让整个 RAG 崩掉，而是降级到另一条路或本地关键词回退。

\subsection{RRF 融合}

RRF 是 Reciprocal Rank Fusion。项目里的 \path{RrfFusionService} 默认 k=60。核心思想是：结果在多个召回列表里排名越靠前，总分越高。

简化公式：

\[
  score(d)=\sum_{s \in S}\frac{1}{k + rank_s(d)}
\]

当前实现还会先归一化各路分数，再把 rank confidence 和 source confidence 组合成最终置信度。这样能把语义检索和关键词检索合成一份统一结果。

\subsection{Rerank 和证据门控}

RAG 不是召回就回答。项目还有两个过滤步骤：

\begin{itemize}
  \item Reranker：如果开启并有可用实现，对候选证据精排。
  \item Evidence Sufficiency Gate：根据知识库 evidenceThreshold 判断证据是否足够。
\end{itemize}

证据不足时，系统返回不足提示，而不是强行编答案。这是项目的核心价值之一。

\section{RAG 问答链路}

\subsection{ConversationController 到 RagPipeline}

流式问答入口：

\begin{lstlisting}[style=code]
POST /api/v1/conversations/{id}/messages/stream
\end{lstlisting}

调用链：

\begin{lstlisting}[style=code]
ConversationController.streamMessage
  -> ConversationService.streamMessage
  -> save user message
  -> build recent conversation history
  -> RagPipeline.streamQuery
  -> save assistant message on complete
\end{lstlisting}

\path{ConversationService} 会保存用户消息，再用最近 10 条消息构造 conversationHistory，交给 QueryRewriteService 改写查询。流结束后，它把助手完整回答和 citations 写入 message 表。

\subsection{RagPipeline 做什么}

\path{RagPipeline} 是 RAG 编排中心：

\begin{enumerate}
  \item 检查用户是否是知识库成员。
  \item 调用 HybridSearchService 检索证据。
  \item 可选 rerank。
  \item 判断证据是否充分。
  \item 用 EvidencePortfolioSelector 选择证据组合。
  \item 渲染 prompt 模板。
  \item 选择 ChatClient。
  \item 流式返回 token、citations、done 事件。
\end{enumerate}

\subsection{证据组合}

\path{EvidencePortfolioSelector.java} 不是简单取 top N。它会综合：

\begin{itemize}
  \item 置信度。
  \item 查询词覆盖。
  \item 文档多样性。
  \item 与已选证据的冗余度。
  \item 上下文长度预算。
\end{itemize}

这让最终 prompt 不只是高分片段堆叠，而是尽量覆盖不同证据来源。

\section{SSE 流式输出}

后端用 Reactor Flux 返回字符串事件。典型事件类型：

\begin{itemize}
  \item token：增量文本。
  \item citations：引用列表。
  \item done：流结束。
  \item error：错误。
\end{itemize}

前端 \path{ChatView.vue} 通过 store 发起请求，使用 \path{authenticatedFetch} 读取流。这样可以同时支持 token 逐步渲染和 401 后刷新 token。

新手调试流式输出时按这个顺序看：

\begin{enumerate}
  \item 浏览器 Network 是否收到 text/event-stream。
  \item 后端是否进入 ConversationService.streamMessage。
  \item RagPipeline 是否返回 token/citations/done。
  \item 前端 store 是否把 token 追加到当前 assistant 消息。
  \item citations 是否能被解析并显示在右侧证据栏。
\end{enumerate}

\section{数据库模型}

\subsection{核心表}

standalone H2 的核心表：

\begin{longtable}{L{0.28\linewidth}L{0.58\linewidth}}
\toprule
表 & 作用 \\
\midrule
sys\_user & 用户账号、密码、角色、状态 \\
sys\_group & 用户默认组和组织信息 \\
group\_member & 用户和组的关系 \\
knowledge\_base & 知识库配置、证据阈值、切片策略 \\
kb\_member & 用户和知识库的成员关系 \\
document & 文档元数据、入库状态、版本信息 \\
document\_chunk & 文档切片、chunk index、ingestionVersion \\
conversation & 会话、绑定知识库、模型 provider \\
message & 用户和助手消息、引用、工具调用记录 \\
analysis\_result & 文件分析结果 \\
refresh\_token & refresh token hash、过期时间、撤销状态 \\
research\_note & 研究笔记、高亮和标签 \\
\bottomrule
\end{longtable}

PostgreSQL 迁移还包含：

\begin{itemize}
  \item document\_chunk\_embedding：向量数据。
  \item citation\_link：引用关系。
  \item kg\_entity、kg\_relation：知识图谱。
  \item document\_permission、audit\_log：文档权限和审计。
  \item scheduled\_task：定时任务。
\end{itemize}

\subsection{读表技巧}

这个项目读数据库不要从 SQL 开始背字段。按业务对象理解：

\begin{longtable}{L{0.30\linewidth}L{0.58\linewidth}}
\toprule
问题 & 主要字段或表 \\
\midrule
用户能看到什么 & sys\_user、sys\_group、group\_member、kb\_member、document\_permission \\
文档能否被问答使用 & document.ingestion\_status、document.active\_ingestion\_version、document\_chunk.ingestion\_version \\
回答能否溯源 & message.citations、document、document\_chunk \\
分析报告在哪里 & analysis\_result、custom.analysis.report-dir \\
\bottomrule
\end{longtable}

\section{文件分析和报告}

项目除了 RAG，还有文件分析功能。相关入口在 \path{AnalysisController.java}，结果由 \path{AnalysisResultService} 保存。报告导出由 \path{ReportExportService} 完成，支持 Markdown 和 PDF。

报告输出目录来自：

\begin{lstlisting}[style=code]
custom.analysis.report-dir=.analysis-reports
\end{lstlisting}

PDF 导出使用 OpenPDF。中文字体优先使用配置的 pdf-font-path，找不到时尝试系统字体，最后回退到内置中文字体。

\section{前后端 API 对照}

\begin{longtable}{L{0.30\linewidth}L{0.58\linewidth}}
\toprule
模块 & 主要接口 \\
\midrule
Auth & \path{POST /api/v1/auth/register}，\path{POST /api/v1/auth/login}，\path{POST /api/v1/auth/refresh}，\path{GET /api/v1/auth/me} \\
Knowledge Base & \path{GET/POST /api/v1/knowledge-bases}，\path{POST /api/v1/knowledge-bases/{id}/search}，成员管理接口 \\
Documents & \path{POST /api/v1/documents/upload}，列表、详情、chunks、retry、批量删除、批量重入库 \\
Conversations & \path{POST /api/v1/conversations}，\path{POST /api/v1/conversations/{id}/messages/stream}，导出、重命名、删除 \\
Citations & \path{POST /api/v1/citations/export} \\
Notes & \path{GET/POST/PUT/DELETE /api/v1/notes} \\
Analysis & \path{/api/analysis/batch}，\path{/api/analysis/results}，Markdown/PDF 导出 \\
Academic & 引用图、引用统计、文献综述 \\
Graph & 知识图谱、实体邻居、路径查询 \\
\bottomrule
\end{longtable}

注意：有些旧接口在 \path{/api} 下，例如 \path{/api/chat/stream} 和 \path{/api/models}；主业务新接口大多在 \path{/api/v1} 下。

\section{安全设计和剩余风险}

\subsection{已具备的安全点}

\begin{itemize}
  \item JWT stateless 认证。
  \item refresh token hash 入库并轮换。
  \item GroupContext 贯穿业务权限。
  \item 文档上传文件名清洗、格式白名单、MIME 检查、大小限制。
  \item 前端请求自动带 token，401 自动刷新。
  \item SecurityConfig 配置 CSP、CORS、无状态 session。
  \item AI 工具的 safePath 限制在 custom.data.base-dir 下，防止路径穿越。
\end{itemize}

\subsection{需要继续加固的点}

这些不是教程推测，而是当前代码和文档反映出的工程边界：

\begin{enumerate}
  \item standalone H2 schema 覆盖核心表，未覆盖 PostgreSQL 迁移的全部能力；完整测试要用 Docker 栈。
  \item \path{GlobalExceptionHandler.handleRuntime} 已返回通用错误信息且不记录异常原文；新增异常处理器时仍要避免把底层 message 返回给客户端。
  \item Swagger、H2 Console、详细运行诊断和 actuator 指标已限制为管理员访问；H2 Console 仅在本地 profile 开启，生产仍需使用受控密钥和网络边界。
  \item \path{AiDataTools} 已定义函数 Bean，但当前 ChatClient 绑定关系需要补验证；不要直接宣称 agent 工具已完整接通。
  \item Docker Compose 要求显式提供敏感凭据，且默认仅暴露前端端口；它用于完整本地集成验证，不等同于生产部署。
\end{enumerate}

\section{部署理解}

\subsection{Standalone}

Standalone 的目标是快：少依赖，适合演示、开发和离线交付。它牺牲了完整生产能力，尤其是 pgvector、Elasticsearch、MinIO、Flyway 迁移链。

适合：

\begin{itemize}
  \item 新手首次启动。
  \item UI 演示。
  \item 不依赖向量检索的基本功能验证。
\end{itemize}

\subsection{Docker Compose}

Docker Compose 的目标是完整：PostgreSQL/pgvector、Elasticsearch、MinIO 和后端一起跑。它适合验证完整 RAG 栈、向量检索、对象存储和 Flyway 迁移。

适合：

\begin{itemize}
  \item 接近生产的功能验证。
  \item 检查集成测试问题。
  \item 验证 PostgreSQL 与 Elasticsearch 的真实行为。
\end{itemize}

\section{工程验证}

本机无 Docker 时，后端本地门禁建议：

\begin{lstlisting}[style=code]
.\mvnw.cmd -q -DskipITs verify
\end{lstlisting}

完整门禁在 Docker 可用机器上跑：

\begin{lstlisting}[style=code]
.\mvnw.cmd verify
\end{lstlisting}

前端验证：

\begin{lstlisting}[style=code]
cd frontend
npm run build
\end{lstlisting}

常见问题：

\begin{longtable}{L{0.30\linewidth}L{0.58\linewidth}}
\toprule
现象 & 判断方式 \\
\midrule
JAVA\_HOME not found & 设置 JDK 21+，本机 README 示例使用 JDK 22 \\
full verify 报 Docker/Testcontainers & 当前机器没有 Docker 或 Docker 不可用，不等于业务代码必然错 \\
前端 401 循环 & 检查 refresh token 是否过期、localStorage 是否有旧 token \\
上传后一直 FAILED & 看 document.failed\_stage 和 error\_message \\
问答没有引用 & 检查文档是否 COMPLETED，检索是否返回结果，证据阈值是否过高 \\
AI 不回答或模型不可用 & 检查 DEEPSEEK\_API\_KEY 或其他 provider key \\
\bottomrule
\end{longtable}

\section{新手二次开发路线}

\subsection{加一个普通页面}

路线：

\begin{enumerate}
  \item 在 \path{frontend/src/views} 新建页面。
  \item 在 \path{frontend/src/router/index.ts} 注册路由。
  \item 在 \path{frontend/src/api} 新增 API 封装。
  \item 后端新增 Controller 和 Service。
  \item 如果需要数据表，加 Flyway migration，同时考虑 H2 standalone 是否要补 schema。
\end{enumerate}

\subsection{加一个新文档格式}

路线：

\begin{enumerate}
  \item 在 DocumentService 的 ALLOWED\_FORMATS 和 MIME 表里加入扩展名。
  \item 在 extractor 包里实现或扩展 FileContentExtractor。
  \item 补一个最小文件解析测试。
  \item 上传样例文件，确认 ETL 到 COMPLETED。
\end{enumerate}

\subsection{加一个检索策略}

路线：

\begin{enumerate}
  \item 新增 SearchService 或扩展已有检索服务。
  \item 在 HybridSearchService 里加入召回结果。
  \item 修改 RRF 融合输入或新增融合策略。
  \item 给降级路径加测试，不要让某一路失败拖垮整个问答。
\end{enumerate}

\subsection{加一个大模型 provider}

路线：

\begin{enumerate}
  \item 在 application.yml 的 custom.ai.providers 下增加 provider。
  \item 确认它兼容 OpenAI API 格式；不兼容就需要适配请求或响应。
  \item 在前端模型选择中显示 provider。
  \item 用 /api/models 检查配置是否读到。
  \item 用 ChatView 发起一次流式问答。
\end{enumerate}

\section{如何调试这个项目}

\subsection{调试上传}

按顺序看：

\begin{enumerate}
  \item 前端上传请求是否命中 \path{/api/v1/documents/upload}。
  \item 后端是否通过 kb\_member 检查。
  \item 文件格式和 MIME 是否通过。
  \item document 是否插入，状态是否 UPLOADED。
  \item EtlPipeline 是否被异步触发。
  \item document.failed\_stage 是否记录失败阶段。
\end{enumerate}

\subsection{调试问答}

按顺序看：

\begin{enumerate}
  \item 会话是否绑定 knowledgeBaseId。
  \item 当前用户是否是知识库成员。
  \item HybridSearchService 是否返回候选结果。
  \item Reranker 是否异常降级。
  \item evidenceThreshold 是否挡住回答。
  \item ChatClient 是否存在。
  \item SSE 是否返回 token、citations、done。
\end{enumerate}

\subsection{调试前端}

按顺序看：

\begin{enumerate}
  \item router 守卫是否因为 token 缺失跳转。
  \item request.ts 是否带 Authorization header。
  \item 401 后 refresh 是否成功。
  \item store 状态是否更新。
  \item 组件 computed 是否从 store 读到数据。
\end{enumerate}

\section{读完后的掌握标准}

如果能回答下面问题，就说明已经真正理解项目：

\begin{enumerate}
  \item 为什么上传后不是马上可问答，而要等 ingestion\_status 到 COMPLETED？
  \item 为什么检索要同时用 pgvector 和 Elasticsearch？
  \item RRF 解决什么问题？
  \item evidenceThreshold 过高会出现什么现象？
  \item GroupContext 从哪里来，为什么 Service 层还要查 kb\_member？
  \item standalone 和 Docker Compose 的数据库能力有什么差异？
  \item 前端为什么同时有 Axios 和 authenticatedFetch？
  \item 如果要加一种文件格式，至少要改哪些地方？
  \item 如果 full verify 失败，如何判断是代码问题还是 Docker/Testcontainers 环境问题？
  \item 当前 agent function 定义和 ChatClient 绑定之间还需要检查什么？
\end{enumerate}

\section{结语}

EviMind 的核心不是“套一个大模型接口”，而是把文档、权限、检索、证据、引用、会话、分析和部署放进一个可运行系统。初学者理解这个项目，要少背框架名，多追数据流：谁创建数据，谁改变状态，谁读取证据，谁负责兜底。把上传链路和问答链路看懂后，其他模块都能沿着同一套思路展开。

\end{document}
